\section{Hierarchical Human-Inspired Memory Architectures}
\label{sec:hierarchical}

This section constitutes the theoretical and empirical core of this survey, examining how principles from human cognitive science have informed the design of memory systems for large language models. The intellectual trajectory from foundational memory theories in experimental psychology, through computational neuroscience models of consolidation and replay, to modern LLM-based agent architectures reveals a striking convergence: the most promising approaches to context management draw explicitly on the hierarchical, multi-store, and consolidation-based architectures that characterize biological memory. We trace this lineage in detail, arguing that the field of LLM context management is, at its conceptual foundations, a branch of applied cognitive science. The section proceeds from classical memory theory through modern implementations, offering both theoretical depth and practical analysis of the systems that operationalize these ideas.

\subsection{Foundations of Human Memory Theory}

The design of memory systems for large language models does not occur in a theoretical vacuum. Over a century of research in experimental psychology has produced detailed models of how biological memory operates, and these models provide the conceptual vocabulary, architectural blueprints, and empirical benchmarks against which artificial memory systems are increasingly evaluated. This subsection reviews the foundational theories that have most directly influenced the design of LLM memory architectures, beginning with the multi-store model of \citet{atkinson1968human} and proceeding through the working memory framework of \citet{baddeley1974working}, the capacity constraints identified by \citet{miller1956magical}, serial position phenomena, the episodic-semantic distinction of \citet{tulving1972episodic}, and the forgetting curve of \citet{ebbinghaus1885memory}.

\subsubsection{The Atkinson-Shiffrin Multi-Store Model}

The multi-store model proposed by \citet{atkinson1968human} established the foundational architecture for understanding human memory as a system of distinct, interacting stores. This model, which has shaped memory research for over fifty years \citep{neath2019fifty}, posits three sequential stages through which information flows. The first stage, sensory memory, functions as a large-capacity buffer that maintains raw sensory information for very brief durations, on the order of hundreds of milliseconds to approximately one second. Sensory memory acts as an input gate, holding a high-fidelity representation of the perceptual world before selective attention determines which information will be passed to the next stage. The second stage, short-term memory (STM), is a severely limited store that holds approximately $7 \pm 2$ chunks of information \citep{miller1956magical} for durations of roughly twenty to thirty seconds without active rehearsal. STM is maintained through control processes such as subvocal rehearsal, and its contents are readily accessible to conscious awareness. The third stage, long-term memory (LTM), is a very large capacity store with potentially lifetime duration, using primarily semantic encoding and susceptible to interference during retrieval rather than simple temporal decay.

Critically, and in a point that is often overlooked in simplified textbook treatments, Atkinson and Shiffrin placed considerable emphasis on the \emph{control processes} that regulate the flow of information across these stores. Attention governs the selection of information from sensory memory into STM. Rehearsal maintains information in STM and facilitates its transfer to LTM. Search and retrieval processes access information stored in LTM, and coding and organizational strategies determine how information is encoded for long-term storage. This emphasis on dynamic processing, rather than merely static stores, is profoundly relevant to modern LLM memory management. The challenge confronting contemporary systems is not storage capacity per se---external databases can store essentially unlimited information---but rather the intelligent movement, selection, and retrieval of information across different storage tiers. The control processes that Atkinson and Shiffrin identified as central to human memory are precisely the mechanisms that LLM memory architectures must implement.

The model maps naturally onto the architecture of modern LLMs. Sensory memory corresponds to token embeddings, which provide a transient, high-dimensional representation of the raw input before any processing by the transformer layers. Short-term memory corresponds to the context window, which functions as a limited-capacity buffer holding the information currently available to the model's attention mechanism. Long-term memory corresponds to the model's parametric knowledge, stored in billions of weight parameters, together with any external databases or retrieval systems that provide persistent storage beyond the context window. The control processes find their analogue in the attention mechanism itself, which dynamically selects and weights information from across the context, and in the various memory management strategies---summarization, retrieval, eviction---that govern what information enters and exits the context window.

The historical impact of the Atkinson-Shiffrin model extends well beyond cognitive psychology. It generated extensive empirical research on memory stores and transitions, directly influenced the development of Baddeley's working memory model, and has served as a conceptual foundation for the design of memory-augmented computational architectures from the earliest expert systems to the most recent LLM-based agents. The model's limitations are also instructive: its assumption of strictly linear, sequential information flow oversimplifies the parallel processing capabilities of biological memory, its treatment of STM as a unitary store was challenged by subsequent research demonstrating domain-specific storage, and its discrete ``stores'' concept has been questioned by neuroscience evidence suggesting more graded representations. These limitations, however, have motivated refinements rather than abandonment, and the core insight---that effective memory requires multiple stores with distinct temporal and capacity characteristics, coordinated by active control processes---remains as relevant to LLM design as it was to cognitive psychology.

\subsubsection{Baddeley's Working Memory Model}

\citet{baddeley1974working} refined and substantially extended the concept of short-term memory by proposing a multi-component \emph{working memory} system. This model challenged the unitary STM view implicit in the Atkinson-Shiffrin framework, replacing it with a more nuanced architecture comprising specialized subsystems coordinated by a central executive. The original 1974 model specified three components, and a fourth was added by \citet{baddeley2000episodic} to address limitations in the original formulation.

\paragraph{The Central Executive.} The central executive is a supervisory attentional control system responsible for managing cognitive resources, coordinating the subsidiary storage systems, and implementing flexible task switching and conflict resolution. It is domain-general and has severely limited capacity, estimated at approximately 1.5 to 2 chunks of actively maintained information. The central executive does not store information itself but rather directs the flow of information between the subsidiary systems and between working memory and long-term memory. It is responsible for the highest-level cognitive operations: selecting strategies, inhibiting irrelevant information, switching between tasks, and allocating processing resources. The underspecification of the central executive---precisely what mechanisms drive its coordination function---has been a persistent criticism of the model \citep{baddeley2003working}, but its conceptual role as an attentional controller remains essential to understanding how limited-capacity systems manage complex information processing.

\paragraph{The Phonological Loop.} The phonological loop is a specialized system for the storage and maintenance of verbal and auditory information. It comprises two subcomponents: a passive phonological store that holds acoustic memory traces for approximately two seconds before they decay, and an active articulatory rehearsal process that refreshes the contents of the store through subvocal repetition. The phonological loop has a capacity of approximately three to four seconds of speech, corresponding to roughly seven to nine words, and demonstrates characteristic phenomena such as the phonological similarity effect (confusions between similar-sounding items) and the word length effect (shorter words recalled better because more can be rehearsed in the two-second window). This component supports language learning, reading comprehension, and verbal problem-solving, and its specialization for sequential linguistic information makes it a particularly apt analogue for the token-level processing that characterizes LLMs.

\paragraph{The Visuospatial Sketchpad.} The visuospatial sketchpad is a system for maintaining and manipulating visual and spatial information. It supports mental rotation, visual problem-solving, navigation, and the maintenance of spatial layouts in working memory. Its capacity is approximately four to five items, and it can be selectively disrupted by concurrent visual tasks without affecting phonological processing, providing strong evidence for the domain-specific separation of working memory components. While LLMs are primarily text-based systems, the visuospatial sketchpad has analogues in the structured representations that LLMs maintain for spatial, relational, and graph-based information within their embedding spaces.

\paragraph{The Episodic Buffer.} Added by \citet{baddeley2000episodic} to address the binding problem---how information from different domains and from long-term memory is integrated into coherent episodes---the episodic buffer is a limited-capacity store holding approximately four chunks of information. It integrates information across domains, binding verbal, visual, and spatial representations into coherent episodes. The episodic buffer serves as the interface between working memory and long-term memory, enabling temporal sequencing of events and episodic narrative construction. Baddeley described it as acting ``as a buffer store, not only between the components of Working Memory, but also linking Working Memory to perception and Long-Term Memory.'' This integrative function is critical: it explains how humans construct coherent representations of ongoing experience from diverse information sources, a challenge that LLMs face when they must synthesize information from different parts of a long context or from external memory systems.

Baddeley's model provides a considerably richer framework for LLM design than the simpler Atkinson-Shiffrin model. The central executive maps onto the transformer's self-attention mechanism, which selects and weights relevant information from across the context and coordinates the processing of different aspects of the input. The phonological loop parallels the token sequence buffer with positional encoding, which maintains sequential linguistic information in a format optimized for language processing. The visuospatial sketchpad corresponds to the multi-dimensional embedding representations that capture structured and relational information. The episodic buffer corresponds to the integrated token representations at the output of transformer layers, which capture multi-modal context bound into coherent representations. Most importantly, the model's emphasis on \emph{active manipulation} of information---not merely passive storage---aligns with the emerging understanding that LLMs must actively manage their context rather than simply accommodating larger inputs. The distinction between storage and processing, which is central to the working memory framework, suggests that context window expansion alone is insufficient; what matters is the quality of the control processes that operate over whatever context is available.

The robust empirical support for the working memory model across diverse experimental paradigms, including demonstrations of domain-specific interference patterns, correlations between working memory capacity and academic achievement, and neuroimaging evidence for component localization, lends confidence to its use as a design framework for artificial systems. The model's limitations---particularly the underspecification of the central executive, the difficulty of measuring capacity consistently across components, and the unclear distinction between refresh and rehearsal mechanisms---are precisely the areas where computational implementations can contribute to theoretical progress by requiring explicit specification of mechanisms that remain implicit in the psychological theory.

\subsubsection{Miller's Magical Number Seven and Chunking}

\citet{miller1956magical} established one of the most enduring findings in cognitive psychology: that the capacity of immediate memory is limited to approximately $7 \pm 2$ ``chunks'' of information, where a chunk is defined as the largest meaningful unit recognized by the individual. This finding has profound implications for information processing because it places an absolute constraint on the amount of information that can be simultaneously maintained and manipulated in working memory. Modern research by \citet{cowan2001magical} has suggested a lower bound of $4 \pm 1$ chunks for the active maintenance capacity of young adults, with Miller's larger estimate potentially reflecting contributions from long-term memory and chunking strategies rather than pure working memory capacity.

The concept of chunking itself---the reorganization of information into larger meaningful units to overcome capacity limitations---is directly relevant to context compression techniques in LLMs. When a chess master perceives a board configuration as a single chunk rather than as thirty-two individual pieces, they dramatically increase the effective information content of their working memory without increasing its capacity in terms of chunks. Analogously, when LLM systems apply gisting, prompt compression, or summarization techniques, they are performing a computational version of chunking: reorganizing detailed information into more compact representations that preserve essential semantic content while reducing the number of tokens consumed in the context window. The effectiveness of these techniques depends on the quality of the chunking---just as expert chess players chunk meaningfully while novices cannot, the quality of LLM compression depends on whether the summarization preserves the information that will be relevant to downstream processing.

The empirical observation that LLMs exhibit working memory-like capacity limits that are surprisingly consistent with human performance---research has shown that GPT-3 and GPT-4 can effectively hold approximately four to seven distinct items in ``active'' consideration, with attention scores gradually concentrating on a small number of positions---suggests that the mathematical structure of attention may impose constraints that are functionally analogous to biological working memory capacity, even though the underlying mechanisms differ fundamentally. This convergence between biological and artificial capacity constraints is theoretically significant because it implies that working memory limitations may reflect computational principles rather than biological accidents, and that strategies for managing limited capacity in human cognition may transfer directly to the design of LLM memory management systems.

\subsubsection{Serial Position Effects: Primacy, Recency, and the Lost-in-the-Middle Phenomenon}

The serial position curve, first characterized systematically by \citet{murdock1962serial} and elaborated by \citet{glanzer1966two}, demonstrates a robust pattern of differential recall based on an item's position within a sequence. Items presented at the beginning of a list are recalled with enhanced accuracy (the primacy effect, attributed to greater rehearsal opportunity and consequent LTM encoding), and items presented at the end of a list are similarly advantaged (the recency effect, attributed to their continued availability in STM at the time of recall). Items in the middle of the list suffer the worst recall, producing the characteristic U-shaped serial position curve. The dissociation between primacy and recency effects---recency is eliminated by a brief period of distractor activity while primacy remains stable---provides strong evidence for the multi-store distinction between STM and LTM.

This phenomenon has a striking and empirically well-documented parallel in the behavior of large language models. \citet{liu2024lost} demonstrated that LLMs exhibit a U-shaped performance curve strikingly similar to the human serial position effect: information placed at the beginning and end of the context window is processed more effectively than information in the middle, a phenomenon they termed ``lost in the middle.'' The cognitive explanation for the human serial position effect---that boundary items receive preferential processing through distinct encoding mechanisms---suggests that the LLM analogue may arise from analogous computational properties of the attention mechanism. Positional embeddings, attention distribution patterns, and the structure of gradient flow through transformer layers may all contribute to preferential processing of boundary positions.

The parallel between human serial position effects and LLM position bias is more than a superficial analogy. In both systems, the effect reflects fundamental properties of how sequential information is processed in a limited-capacity system. The strategies that humans deploy to overcome serial position bias---such as rehearsal, elaborative encoding, and organizational strategies that reduce reliance on position---may inform computational strategies for mitigating the lost-in-the-middle effect in LLMs, such as strategic information placement, iterative re-reading, and architectural modifications that equalize attention across positions.

\subsubsection{Tulving's Episodic and Semantic Memory Distinction}

\citet{tulving1972episodic} introduced one of the most influential distinctions in memory research: the separation of long-term memory into \emph{episodic memory} for specific events grounded in time and place, and \emph{semantic memory} for general knowledge, facts, and language. This distinction was further refined by \citet{tulving1985memory}, who associated each memory type with a distinct form of consciousness: episodic memory involves autonoetic consciousness (the ability to mentally re-experience past events, to ``travel back in time''), semantic memory involves noetic consciousness (abstract knowing without re-experiencing), and procedural memory involves anoetic consciousness (non-conscious skilled performance).

The episodic-semantic distinction maps directly onto the architecture of LLM memory systems. Parametric knowledge acquired during pre-training represents semantic memory: general truths, factual relationships, and linguistic patterns stored implicitly in the model's weight parameters. This knowledge is accessed through the forward pass without any explicit retrieval mechanism, analogous to how humans access semantic knowledge without consciously recalling when or where they learned it. Context-specific information from particular interactions---conversation histories, user preferences, specific task details---represents episodic memory: particular events linked to temporal context and situational specifics.

Standard LLMs excel at semantic knowledge but struggle profoundly with episodic memory. They cannot, without external augmentation, remember specific past interactions, track the temporal context of information, or maintain spatiotemporal grounding across sessions. This limitation has been identified as the ``missing piece'' for long-term LLM agents \citep{sun2025episodic}. The challenges are multifaceted: difficulty with information presented only once over long timescales, poor performance on sequence order recall tasks, interference from multiple related events stored in the same parametric space, and the fundamental problem that parametric knowledge is updated only through training, not through experience \citep{ning2025episodic}. Solutions to these challenges---external episodic memory systems such as EM-LLM \citep{hu2025emllm}, explicit temporal tagging and event segmentation, and structured retrieval combining recency, relevance, and importance---all draw directly on Tulving's theoretical framework, implementing computational versions of the episodic memory system that standard LLMs lack.

Modern cognitive science has increasingly viewed episodic and semantic memory as existing on a continuum rather than as rigidly discrete categories, with shared underlying neural processes and gradual transitions from episodic to semantic representations through consolidation. This perspective is equally relevant to LLM design: the most effective memory architectures may be those that support gradual transformation of specific episodic information into generalized semantic knowledge, mirroring the consolidation processes that operate in biological memory.

\subsubsection{The Ebbinghaus Forgetting Curve}

\citet{ebbinghaus1885memory} conducted the first systematic experimental investigation of memory, establishing through heroic self-experimentation that memory retention follows an exponential decay function. The mathematical form of the forgetting curve is typically expressed as:
\begin{equation}
R = e^{-t/S}
\end{equation}
where $R$ represents retention (as a proportion of original learning), $t$ is the time elapsed since learning, and $S$ is a parameter representing memory strength that varies with the depth of encoding, the number of repetitions, and individual differences. The forgetting curve exhibits a characteristic shape: a steep initial decline, with approximately fifty to seventy percent of information lost within the first twenty-four hours without reinforcement, followed by a gradual leveling off at longer intervals. This pattern was replicated and confirmed by \citet{murre2015replication} using modern experimental methods, validating Ebbinghaus's original findings more than a century after their initial publication.

The debate between exponential and power-law decay functions continues in the memory literature, with evidence suggesting that pure forgetting of homogeneous material follows exponential decay while the aggregate forgetting of heterogeneous material---where different items have different memory strengths---may produce power-law curves that are the superposition of multiple exponential functions. This distinction is relevant to LLM memory design because the information stored in artificial memory systems is invariably heterogeneous, suggesting that power-law decay models may be more appropriate than simple exponential decay for practical implementations.

Ebbinghaus also discovered the spacing effect---that distributed practice produces more durable learning than massed practice---and the related finding that retrieval practice (testing) produces stronger memory than passive review. These findings have been operationalized in spaced repetition algorithms that schedule reviews to maximize retention with minimum study time. The implications for LLM memory management are direct: information that is periodically retrieved and reinforced should be maintained more durably than information that is merely stored, and the scheduling of memory access and consolidation operations should follow principles analogous to spaced repetition.

The Ebbinghaus forgetting curve has been directly implemented in LLM memory systems. \citet{zhong2024memorybank}, in their MemoryBank system, implemented exponential decay of memory importance scores to determine which stored information should be retained versus allowed to fade, creating an artificial analogue of biological forgetting. This approach represents a departure from most LLM memory systems, which either retain everything indefinitely or use attention-based selection without temporal decay, and it highlights the insight that principled forgetting---not merely storage---is essential for scalable and effective memory management.

\subsection{Memory Consolidation and Complementary Learning Systems}

The transition of information from short-term to long-term storage is not a simple copying process but involves active transformation through consolidation. This subsection examines the neurobiological theories of memory consolidation that have most directly influenced the design of LLM memory architectures, focusing on the complementary learning systems framework, hippocampal replay, and sleep-dependent consolidation.

\subsubsection{Systems Consolidation: From Hippocampus to Neocortex}

Memory consolidation operates on two distinct timescales. Synaptic consolidation, occurring over minutes to hours, involves molecular cascades that strengthen synaptic connections at the cellular level through mechanisms including AMPA receptor insertion, protein synthesis, and calcium-dependent signaling. Systems consolidation, occurring over hours to years, involves a large-scale reorganization in which memories initially dependent on the hippocampus are gradually redistributed to neocortical networks, becoming increasingly hippocampus-independent over time. The empirical foundation for this distinction was established through lesion studies, most notably the case of patient H.M., whose selective hippocampal damage preserved remote memories while severely impairing new episodic learning \citep{squire1991medial}. This dissociation demonstrated that the hippocampus is essential for the formation of new declarative memories but not for the storage of old ones, implying a process by which memories are transferred from hippocampal to neocortical substrates.

The distinction between synaptic and systems consolidation has a direct analogue in LLM memory design. Synaptic consolidation corresponds to the immediate processing that occurs when new information enters the context window---the rapid formation of representations through attention and feedforward computation. Systems consolidation corresponds to the longer-term processes by which information from the context window is selectively transferred to external memory stores, summarized, and integrated with existing knowledge. The absence of principled systems consolidation in standard LLMs---information either remains in the context window or is lost entirely when the window advances---represents a fundamental architectural limitation that motivates the development of explicit consolidation mechanisms in memory-augmented systems.

\subsubsection{Complementary Learning Systems Theory}

\citet{mcclelland1995complementary} proposed the complementary learning systems (CLS) framework to explain a puzzle that had long troubled connectionist modelers: why does the brain require two specialized learning systems rather than a single, unified network? The answer, McClelland and colleagues argued, lies in a fundamental tension between the requirements of rapid learning and the requirements of structured knowledge representation.

The hippocampal system implements rapid, one-shot learning through sparse, pattern-separated representations. Its encoding is pattern-specific and detail-preserving, allowing new experiences to be stored in a single exposure without disrupting existing memories. However, the hippocampus has limited capacity and stores memories that are relatively short-lived, on the order of days to weeks. The neocortical system, by contrast, implements gradual, interleaved learning through distributed, overlapping representations. It builds abstract, generalizable knowledge structures over many exposures, but it cannot rapidly incorporate new information without risking catastrophic interference with existing representations.

The consolidation process bridges these two systems. Memories are first rapidly encoded in the hippocampus, then gradually transferred to the neocortex through repeated reinstatement---a process in which hippocampal memory traces are replayed to the neocortex, allowing the neocortex to learn from these replayed experiences at a rate that avoids catastrophic interference. This interleaved replay ensures that new memories are integrated with existing knowledge structures rather than overwriting them.

The CLS framework provides a biological blueprint for LLM memory design that has been enormously influential. Current LLMs function primarily as neocortical systems: they build distributed knowledge representations through gradual training on large datasets, but they lack an equivalent of the hippocampal system that would allow rapid encoding of new episodic information without disrupting existing parametric knowledge. Memory-augmented architectures that combine rapid external storage for episodic information with stable parametric knowledge for semantic information effectively implement a computational analogue of complementary learning systems. The external memory store (whether a vector database, a knowledge graph, or a structured memory bank) functions as the hippocampal system, providing rapid, pattern-separated storage for new information. The model's parameters function as the neocortical system, providing stable, distributed storage for generalized knowledge. The retrieval and integration mechanisms that connect these two stores correspond to the consolidation processes that transfer information between hippocampus and neocortex.

\paragraph{Pattern Separation and Pattern Completion.} Two key computational principles from the CLS framework are directly applicable to LLM memory design. Pattern separation---the creation of distinct, non-overlapping representations for similar but different memories---is essential for preventing interference between stored memories. In the hippocampus, this is achieved through sparse coding in the dentate gyrus. In LLM memory systems, analogous mechanisms include the use of distinct embedding spaces for different memory types, the explicit separation of episodic and semantic memory stores, and the use of hash-based or graph-based indexing that maintains separation between similar but distinct memories. Pattern completion---the retrieval of a full memory from a partial cue---is the complementary process that enables effective memory retrieval. In the hippocampus, pattern completion operates through recurrent connections in CA3 that allow partial activation patterns to recover full stored patterns. In LLM memory systems, this corresponds to similarity-based retrieval from vector databases, where a query embedding is matched against stored memory embeddings to recover the full content of the most similar memories.

\subsubsection{Hippocampal Replay and Sleep-Dependent Consolidation}

Neural replay mechanisms support memory consolidation through the sequential reactivation of experience patterns during both waking rest and sleep \citep{carr2011hippocampal, foster2017replay}. During waking rest periods, the hippocampus frequently reactivates recent experience patterns, a process that may serve multiple functions including consolidation, planning, and memory retrieval. During sleep, replay becomes particularly organized and effective, occurring within a framework of coordinated oscillatory activity that provides optimal conditions for hippocampal-cortical information transfer.

Sleep-dependent consolidation operates through several complementary mechanisms \citep{diekelmann2010memory}. During non-rapid eye movement (NREM) sleep, particularly slow-wave sleep, three types of oscillatory activity coordinate memory consolidation. Slow oscillations at approximately 0.5 to 1 Hz coordinate large-scale information transfer between hippocampus and neocortex. Sleep spindles at 12 to 16 Hz, generated by thalamic circuits, mark periods of active memory replay and facilitate synaptic plasticity in cortical networks. Sharp-wave ripples at 140 to 200 Hz, generated in the hippocampus, represent the actual reactivation of specific memory traces. The temporal coupling of these oscillations---ripples nested within spindles, which are in turn nested within slow oscillations---creates a hierarchical timing structure that orchestrates the precise sequence of events required for effective consolidation. Research has demonstrated that sleep enhances memory consolidation by twenty-five to forty percent compared to equivalent periods of wakefulness, and that targeted memory reactivation during sleep (presenting cues associated with specific memories) can selectively strengthen particular memories \citep{rasch2013sleep}.

During rapid eye movement (REM) sleep, a distinct set of mechanisms operates. Theta-frequency oscillations at 4 to 8 Hz dominate hippocampal activity, and reduced norepinephrine levels enhance synaptic plasticity. REM sleep appears to play a complementary role to NREM sleep in consolidation, with particular importance for emotional memory processing, the integration of new memories into existing semantic networks, and the extraction of abstract patterns from specific experiences. The alternation of NREM and REM sleep stages across the night facilitates what \citet{gonzalez2022model} have termed ``graceful integration''---the gradual incorporation of new information into existing knowledge structures without disrupting previously consolidated memories.

These biological mechanisms have no direct analogue in standard LLM training, which relies exclusively on online gradient-based learning from sequentially presented training examples. The absence of offline consolidation phases---analogous to sleep---is arguably a fundamental limitation of current LLM architectures. Standard LLMs have no mechanism for replaying and reorganizing recently acquired information, no process for integrating new knowledge with existing parametric representations without catastrophic interference, and no offline phase during which abstract patterns can be extracted from specific training experiences. This gap has motivated a growing body of research on sleep-inspired training procedures, discussed in detail below, that aim to introduce consolidation-like mechanisms into the training of artificial neural networks.

\subsection{Classical Cognitive Architectures}

Before the advent of large language models, the field of cognitive architecture had already established fundamental principles about how memory systems should be organized for intelligent behavior. Two architectures---SOAR and ACT-R---are particularly relevant because they formalized many of the memory distinctions that have been independently rediscovered in the design of LLM-based agents.

\subsubsection{SOAR}

The SOAR cognitive architecture \citep{laird2012soar} is a production system architecture that models intelligent behavior as the interaction between working memory and procedural knowledge encoded as production rules (if-then statements that match conditions in working memory to actions). SOAR represents working memory as relational graph structures composed of node-edge-value triples, providing a structured representation that captures relationships between entities rather than merely listing features. Procedural memory in SOAR stores the production rules that encode the agent's behavioral repertoire, and the architecture includes mechanisms for learning new productions through experience (chunking) and for resolving impasses when no production matches the current working memory state.

SOAR's contribution to the understanding of memory in intelligent systems lies in its demonstration that effective behavior requires not merely storage but structured interaction between working memory and long-term procedural knowledge. The architecture's focus on the dynamic interplay between current state (working memory) and accumulated expertise (procedural memory) anticipates the design of modern agent architectures, where the LLM's context window serves as working memory and its parametric knowledge serves as a form of procedural-semantic memory. SOAR also demonstrated the importance of learning from experience---its chunking mechanism, which compiles frequently used reasoning sequences into compiled productions, parallels the reflection and abstraction mechanisms found in modern generative agent systems.

\subsubsection{ACT-R}

The ACT-R architecture \citep{anderson2004integrated} provides a more detailed cognitive model than SOAR, with specialized modules for different cognitive functions (visual, memory, motor, goal) coordinated by a central production system. ACT-R's memory architecture distinguishes between declarative memory (facts and episodic knowledge stored as chunks with associated activation levels) and procedural memory (skills and behavioral routines stored as production rules). Declarative memory retrieval in ACT-R follows a rational analysis of memory: the probability of retrieving a memory chunk is a function of its base-level activation (reflecting recency and frequency of use) and its associative activation (reflecting the degree to which current context provides retrieval cues for the chunk).

The mathematical formulation of ACT-R's memory retrieval---particularly its base-level learning equation, which models memory activation as a function of recency and frequency---provided an early computational specification of the principles that Ebbinghaus identified empirically. The architecture's distinction between declarative and procedural memory directly parallels the semantic and procedural memory distinction in modern language agent frameworks such as CoALA. ACT-R's integration of multiple specialized modules under central executive control also anticipates the modular architectures that characterize many contemporary LLM-based agent systems, where different modules handle different aspects of memory, retrieval, and action.

\paragraph{Lessons from Classical Architectures.} Both SOAR and ACT-R demonstrate a principle that has been rediscovered repeatedly in the design of LLM-based agents: effective intelligent behavior requires multiple memory systems with distinct access patterns, storage durations, and consolidation mechanisms. A single, undifferentiated memory store is insufficient for complex cognitive tasks. The structured interaction between working memory and long-term stores, the distinction between declarative and procedural knowledge, and the importance of learning mechanisms that refine and consolidate knowledge through experience---all principles established by classical cognitive architectures---are now being implemented in modern LLM memory systems, often with explicit acknowledgment of their cognitive science origins.

\subsection{Modern LLM Memory Architectures Inspired by Human Cognition}

The theoretical foundations reviewed above have directly inspired a generation of memory architectures for large language models and LLM-based agents. This subsection examines the most influential of these architectures in detail, analyzing how each one draws on cognitive science principles and evaluating its contributions and limitations.

\subsubsection{CoALA: Cognitive Architectures for Language Agents}

\citet{sumers2024cognitive} proposed CoALA (Cognitive Architectures for Language Agents), the first comprehensive framework that systematically organizes the design space of language agents using concepts drawn from classical cognitive architectures. CoALA addresses a gap in the rapidly growing literature on LLM-based agents: while numerous systems had been proposed, no systematic framework existed to organize and compare them, to identify design tradeoffs, or to connect modern implementations to the decades of research on cognitive architectures that preceded them. CoALA provides this unifying lens by organizing language agents along three fundamental dimensions inspired by cognitive science.

\paragraph{Information Storage.} The first dimension concerns how agents store and organize information. CoALA distinguishes between working memory, which corresponds to the context window and holds the current state, recent observations, and active goals, and long-term memory, which encompasses three subtypes derived from cognitive psychology. Procedural memory stores skills and strategies encoded in model parameters through gradient descent, corresponding to the implicit knowledge that an LLM acquires during pre-training and fine-tuning. This includes learned patterns for tool use, reasoning strategies, and behavioral routines that are executed automatically during inference. Semantic memory stores facts and concepts in learned embeddings and, optionally, in external knowledge bases, corresponding to the explicit factual knowledge that can be retrieved through the forward pass or through external retrieval mechanisms. Episodic memory stores specific events and experiences in external memory banks, corresponding to the temporally grounded, context-specific memories that standard LLMs lack but that are essential for agents operating over extended interactions. The CoALA framework makes explicit a distinction that is often implicit in the design of LLM agents: the context window is working memory, not the entirety of memory, and effective agent design requires consideration of how information flows between working memory and the various forms of long-term storage.

\paragraph{Action Space.} The second dimension concerns the actions available to the agent, divided into external actions that affect the environment (API calls, tool use, environmental interaction) and internal actions that operate on the agent's own representations (reasoning, memory retrieval, and learning). This distinction is important because it highlights that memory operations---storing, retrieving, and updating information---are themselves actions that the agent must decide to perform, not passive processes that occur automatically. The agent must decide when to store information in long-term memory, when to retrieve information from long-term memory into working memory, and when to update or consolidate stored information. These decisions are analogous to the control processes that Atkinson and Shiffrin identified as central to human memory management.

\paragraph{Decision-Making.} The third dimension concerns how the agent decides which action to take, modeled as an interactive loop with planning and execution phases. The planning phase involves perception (observing the current state), goal activation (determining what to pursue), and deliberation (planning the next action). The execution phase involves performing the chosen action, observing the results, and updating working memory. This loop structure mirrors the perception-action cycle of classical cognitive architectures and provides a framework for understanding how memory operations are integrated into the agent's overall behavior.

CoALA's contribution extends beyond taxonomy. By retrospectively organizing over one hundred papers on language agents and mapping each to its framework, CoALA identifies systematic gaps in the literature. Most current systems, the analysis reveals, lack robust episodic memory and principled consolidation mechanisms. Procedural learning---the ability to acquire new skills and strategies through experience---is particularly underexplored. The framework also identifies promising but underexplored design combinations, such as agents with rich episodic memory combined with explicit planning capabilities, or agents with procedural learning mechanisms that can consolidate frequently used reasoning patterns into efficient routines. By connecting modern LLM agents to classical cognitive architectures such as ACT-R, SOAR, and CLARION, CoALA establishes a conceptual bridge that enables the transfer of insights from decades of cognitive science research to the design of contemporary systems.

\subsubsection{Generative Agents: Memory Streams and Reflection}

\citet{park2023generative} developed generative agents---computational agents embedded in a simulated environment that exhibit believable human-like behavior through a hierarchical memory architecture that operationalizes episodic memory for language agents. This work represents a landmark demonstration that human-inspired memory systems can enable emergent complexity in artificial agents, and it has been enormously influential in subsequent research on agent memory design.

\paragraph{The Memory Stream.} The core data structure underlying generative agents is the memory stream: a chronologically ordered database of every observation, action, plan, and reflection that an agent experiences. Each entry in the memory stream is stored as a natural language description accompanied by a timestamp, an importance score (rated on a scale of 0 to 10 by the language model), and a vector embedding that enables similarity-based retrieval. The memory stream grows unboundedly over the course of the simulation, capturing the complete experiential history of the agent in a format that is both human-readable and computationally searchable. This design choice---storing memories as natural language rather than as structured data or compressed representations---reflects a commitment to interpretability and flexibility, allowing the same retrieval mechanisms to operate over observations, actions, and abstract reflections.

\paragraph{The Reflection Mechanism.} The reflection mechanism is the key innovation that transforms the memory stream from a passive log into an active knowledge construction system. Periodically, when the cumulative importance scores of recent unreflected memories exceed a threshold, the agent generates higher-level abstractions from its recent experiences. The reflection process involves querying the language model to identify patterns, draw conclusions, and form beliefs based on a collection of recent memories. The resulting reflections are themselves stored as entries in the memory stream, creating a hierarchical structure: raw observations at the leaf level, event-level summaries at intermediate levels, and abstract beliefs and goals at the highest level. Because reflections are stored alongside observations and can themselves be reflected upon, the system supports meta-reflection---the formation of increasingly abstract and general knowledge from specific experiences.

In practice, agents in the Smallville simulation reflected two to three times per simulated day. The reflection process closely parallels the consolidation mechanisms described in the CLS framework: specific episodic experiences are gradually abstracted into generalized semantic knowledge through a process of repeated review and integration. The analogy to sleep-dependent consolidation is also apparent: reflection represents an offline processing phase during which the agent steps back from active engagement with the environment to reorganize and abstract from its accumulated experience.

\paragraph{Retrieval Mechanism.} When planning actions, agents retrieve relevant memories from the memory stream using a weighted combination of three signals:
\begin{equation}
\text{score}(m) = \alpha \cdot \text{recency}(m) + \beta \cdot \text{importance}(m) + \gamma \cdot \text{relevance}(m, q)
\end{equation}
where recency decays exponentially with time since the memory was last accessed, importance reflects the rated significance of the memory, and relevance is computed as the cosine similarity between the memory's embedding and the current query embedding. This multi-signal retrieval mechanism draws on multiple principles from cognitive psychology: recency-weighted retrieval parallels the recency effect and the Ebbinghaus forgetting curve, importance-weighted retrieval parallels the preferential encoding of emotionally significant or goal-relevant information, and relevance-based retrieval parallels the associative, cue-dependent retrieval that characterizes human episodic memory.

\paragraph{Emergent Behaviors and Evaluation.} In a simulation of twenty-five agents in a virtual village (``Smallville''), powered by GPT-3.5-turbo, the full memory architecture enabled a remarkable range of emergent behaviors that were not explicitly programmed. Agents maintained consistent daily schedules influenced by their accumulated memories. Relationships developed naturally through repeated interactions, with agents greeting friends differently from strangers and conversations reflecting the history of previous interactions. Information spread through social networks in realistic patterns, with gossip propagating believably and at realistic velocities. Most strikingly, agents spontaneously coordinated group activities: a Valentine's Day party was organized through a series of conversations among multiple agents, each independently deciding to attend based on their own memories and social relationships.

Human raters, blind to experimental condition, evaluated agent believability on scales of consistency, appropriateness, and realism. The full memory system (memory stream plus reflection) received an overall believability rating of 7.8 out of 10, compared to 5.2 without reflection and 3.1 without any memory stream. Random action agents received only 1.9. The 2.5-fold improvement from adding reflection alone demonstrates the critical importance of abstraction and consolidation mechanisms---raw episodic storage is necessary but insufficient for believable behavior; the consolidation of specific experiences into generalized knowledge is what enables coherent, adaptive, contextually appropriate behavior over extended periods.

\paragraph{Limitations and Future Directions.} Despite its success, the generative agents architecture has several important limitations. It lacks explicit forgetting mechanisms: the memory stream grows without bound, and no information is ever discarded, leading to increasing computational costs for retrieval as the simulation progresses. There is no sleep-like consolidation phase that would enable offline reorganization of memories. The reflection mechanism operates only on recent memories, potentially missing long-term patterns that span many simulated days. Hallucination is possible, with agents occasionally referencing events that did not occur. And the computational cost is substantial: with twenty-five agents each making multiple LLM calls per simulation step, a full simulation requires approximately twenty-five thousand LLM calls. These limitations define the research agenda for subsequent memory architectures, many of which address specific shortcomings of the generative agents approach.

\subsubsection{MemGPT: Virtual Context Management as an Operating System}

\citet{packer2023memgpt} introduced MemGPT, an architecture that draws inspiration not from cognitive psychology but from operating systems engineering, specifically from the concept of virtual memory. In computer operating systems, virtual memory creates the illusion of a large, uniform address space by intelligently managing the movement of data between fast but limited primary memory (RAM) and slow but capacious secondary storage (disk). MemGPT applies this principle to LLM context management, treating the context window as ``main memory'' and external storage as ``disk,'' with the LLM itself serving as the processor that issues instructions for data movement between tiers.

The key innovation of MemGPT is that the LLM agent manages its own memory through function calls. When the agent determines that it needs information not currently in the context window, it issues a retrieval function call to load relevant information from external storage. When the context window approaches capacity, the agent issues storage function calls to offload less immediately relevant information to external storage. An interrupt mechanism manages control flow when context limits are reached, pausing the current operation to perform memory management before resuming. This self-managed approach contrasts with systems where memory management is handled by external controllers: in MemGPT, the LLM is both the processor and the memory manager, deciding what to page in and out of its own context.

MemGPT demonstrates two primary applications. For document analysis, it enables processing of documents that far exceed the context window by intelligently paging relevant sections into the context as needed. For multi-session chat, it maintains long-term memory across conversation sessions by storing and retrieving conversation summaries, user preferences, and factual information from persistent external storage. The result is an illusion of infinite memory through virtualization, analogous to how virtual memory creates the illusion of unlimited RAM.

While MemGPT's inspiration is primarily computational rather than cognitive, the resulting architecture has clear parallels to hierarchical memory models from cognitive science. The context window functions as working memory with limited capacity and immediate accessibility. External storage functions as long-term memory with unlimited capacity but slower access requiring explicit retrieval operations. The memory management operations---deciding what to store, what to retrieve, and what to evict---correspond to the control processes that Atkinson and Shiffrin identified as central to human memory management. The system's ability to create the illusion of continuous memory across sessions, despite the fundamentally stateless nature of each individual LLM inference, parallels the way that human memory creates the illusion of continuous experience despite the discontinuities introduced by sleep, distraction, and the limited capacity of working memory.

\subsubsection{MemoryBank: Ebbinghaus-Inspired Forgetting}

\citet{zhong2024memorybank} directly implemented the Ebbinghaus forgetting curve in an LLM memory system, creating what is arguably the most explicit translation of classical memory theory into computational practice. MemoryBank stores past conversations, summarized events, and user portraits in a structured memory bank, with each memory assigned an importance score that decays exponentially over time according to a modified Ebbinghaus function:
\begin{equation}
\text{Importance}(m, t) = I_0 \cdot e^{-\lambda t}
\end{equation}
where $I_0$ is the initial importance score, $t$ is the time elapsed since last access, and $\lambda$ is a decay rate controlling the speed of forgetting.

The system comprises three modules that work together to manage the memory lifecycle. The Writer module extracts important information from each dialogue turn---user preferences, emotional states, factual details, goals---and stores it with associated importance scores. The Retriever module finds relevant memories for the current dialogue context using a combination of semantic similarity, recency, and importance scoring. The Reader module integrates retrieved memories into the LLM's input prompt, formatting them as natural language context that augments the model's understanding of the user and the conversation history.

The critical innovation is the memory update mechanism. Memories that are not accessed undergo exponential decay: a memory created with importance 8.0 might decline to 7.1 after one day, 5.9 after three days, 3.2 after a week, and 0.9 after two weeks, at which point it falls below the retrieval threshold and effectively becomes inaccessible. However, when a memory is retrieved and used in dialogue, its importance is boosted, resetting the decay curve. This reinforcement mechanism mirrors the spacing effect: frequently relevant memories are strengthened through repeated retrieval, while irrelevant memories naturally fade. The result is an adaptive memory system in which the contents are shaped by relevance and usage patterns rather than by arbitrary retention or deletion rules.

MemoryBank was demonstrated through SiliconFriend, an AI companion chatbot that leverages the memory system to provide empathetic responses, personality understanding, and the impression of long-term companionship. Users interacting with SiliconFriend reported more personalized and contextually appropriate responses compared to systems without persistent memory. The explicit incorporation of a principled forgetting mechanism represents a significant departure from most LLM memory systems, which either retain everything indefinitely (leading to unbounded growth and retrieval degradation) or use ad hoc deletion strategies without theoretical motivation. MemoryBank demonstrates that forgetting is not merely a limitation to be overcome but an essential feature of effective memory management---a principle that Ebbinghaus established empirically and that modern information theory supports through the concept of optimal information retention under resource constraints.

\subsubsection{Think-in-Memory: Storing Thoughts, Not Facts}

\citet{liu2023thinkinmemory} proposed Think-in-Memory (TiM), an architecture based on the observation that humans recall the results of previous thinking without re-reasoning from scratch, whereas LLMs repeatedly perform recall-reason cycles that produce biased and inconsistent results. When a human encounters a problem similar to one they have solved before, they recall their previous solution or conclusion directly, rather than re-deriving it from first principles. Standard LLMs, lacking persistent memory of their own reasoning, must re-reason every time, and this repeated reasoning is susceptible to inconsistency---the same question posed in different contexts may elicit different reasoning chains and different conclusions.

TiM addresses this problem by storing \emph{post-thinking results}---the conclusions and insights generated by the model's reasoning process---rather than raw input history. The architecture operates in two phases. In the pre-response phase, the system recalls relevant thoughts from memory before generating a response, providing the model with the results of its own previous reasoning as additional context. In the post-response phase, after generating a response, the system extracts the key thoughts and insights from the reasoning process and updates memory with these new thoughts. Memory operations include insertion of new thoughts, forgetting of outdated thoughts, and merging of similar thoughts based on similarity computed via Locality-Sensitive Hashing, which provides efficient approximate nearest-neighbor search for memory management.

The insight underlying TiM---that storing processed thoughts is more efficient and effective than storing raw observations---has a direct analogue in cognitive science. Human long-term memory does not store verbatim records of experience but rather stores the gist, the meaning, and the conclusions derived from experience. The levels-of-processing framework \citep{craik1972levels} established that deeper semantic processing produces more durable memory traces than shallow perceptual processing, a finding that supports TiM's approach of storing the products of deep reasoning rather than the raw inputs to reasoning.

\subsubsection{Self-Controlled Memory}

\citet{liang2023scm} introduced Self-Controlled Memory (SCM), which addresses the specific problem of noise accumulation when external memory is injected into the LLM context. When information from an external memory store is retrieved and inserted into the context window, it may introduce irrelevant or outdated information that degrades rather than improves the model's performance. SCM addresses this through a memory controller that determines when and how to use stored information, implementing a form of metacognitive control over memory access.

SCM employs two memory types per processing segment. Activation memory stores long-term historical information that has been consolidated from previous processing segments, representing the accumulated knowledge relevant to the current task. Flash memory stores short-term, real-time information from the immediately preceding segment, representing the most recent context. The memory controller, implemented as a meta-level decision-making component, evaluates whether retrieved memories are relevant to the current processing context and controls their injection into the model's input. This plug-and-play design requires no fine-tuning and handles ultra-long texts through selective memory injection, achieving seventy-eight percent retrieval recall compared to forty-five percent for baseline approaches on long-context tasks.

The distinction between activation memory and flash memory in SCM parallels the distinction between long-term and short-term memory in the Atkinson-Shiffrin model, while the memory controller's function of evaluating and selectively admitting external information into the processing context parallels the attentional gating function of the central executive in Baddeley's working memory model. The key contribution of SCM is its demonstration that uncontrolled memory injection can be counterproductive, and that intelligent memory management requires not only effective storage and retrieval but also effective filtering and selection at the point of integration.

\subsubsection{RecallM: Graph-Based Temporal Memory}

\citet{kynoch2023recallm} introduced RecallM, which departs from the dominant paradigm of vector-database-based memory by using a neuro-symbolic graph database for relational modeling. The key insight motivating this approach is that vector databases, while effective for semantic similarity search, cannot naturally represent relationships between concepts, temporal ordering of events, or the revision of beliefs over time. A graph database, by contrast, can explicitly represent entities, relationships, and temporal information, enabling more sophisticated reasoning about the structure and dynamics of stored knowledge.

RecallM's key advantages include advanced relationship modeling between concepts, explicit temporal understanding that tracks when information was learned and how beliefs have changed over time, and one-shot belief updating that allows the system to revise prior beliefs based on new evidence. The one-shot belief updating capability is particularly significant: in a graph database, contradictory information can be resolved by updating relationship edges rather than by attempting to reconcile conflicting entries in a vector store. A hybrid version combining graph and vector databases (Hybrid-RecallM) leverages the complementary strengths of relational and semantic representations, using the graph database for structured knowledge and belief management and the vector database for flexible semantic retrieval.

RecallM's approach has connections to the semantic network models of human memory that were influential in cognitive psychology during the 1960s and 1970s, particularly the spreading activation model of \citet{collins1975spreading}. In these models, concepts are represented as nodes in a network, with relationships represented as edges, and memory retrieval operates through a process of spreading activation from cue nodes to target nodes. RecallM's graph database implements a computational version of this spreading activation process, with explicit representation of the relational structure that underlies semantic knowledge.

\subsubsection{RET-LLM: Triplet-Based Knowledge}

\citet{modarressi2024retllm} implemented a read-write memory system for LLMs based on semantic triplets (subject-predicate-object), grounded in Davidsonian event semantics. Davidsonian semantics represents events as logical entities with participants and properties, providing a formal framework for decomposing complex statements into structured, composable units. By storing knowledge as triplets, RET-LLM achieves four properties that distinguish it from opaque vector stores: scalability (the memory grows naturally with new information), aggregatability (related facts can be combined and summarized), updatability (existing knowledge can be revised when new information contradicts it), and interpretability (the stored knowledge is in human-readable form). A memory controller manages the lifecycle of stored knowledge, handling extraction of triplets from text, lookup of relevant triplets given a query, and generation of responses grounded in retrieved factual triplets. RET-LLM achieved superior performance on temporal question answering tasks, where the ability to track when facts were learned and how they have changed is essential for accurate responses.

\subsubsection{RAISE: Dual-Component Memory for Conversational Agents}

\citet{liu2024raise} enhanced conversational agents with a dual-component memory architecture that mirrors the human STM/LTM distinction. Dialogue-level memory maintains the current conversation history and a scratchpad for intermediate reasoning, functioning as a form of working memory that supports ongoing processing. Turn-level memory stores query-response examples and task trajectories from previous interactions, functioning as a form of long-term episodic and procedural memory that enables the agent to draw on past experience when handling new situations. Fine-tuning the agent on domain-specific examples enhances the controllability and efficiency of memory operations, making the approach particularly effective in specialized applications such as real estate sales, customer service, and technical support, where the ability to recall and apply domain-specific knowledge from previous interactions is essential for effective performance.

\subsection{Transformer Attention as Working Memory Analogue}

The relationship between transformer attention and human working memory extends beyond loose analogy to reveal deep structural and functional parallels that illuminate both biological and artificial cognition.

\subsubsection{Cue-Based Retrieval and Attention}

Transformer attention implements a form of cue-based retrieval that closely parallels the associative retrieval mechanisms described in cognitive models of human memory. In the attention mechanism, queries serve as retrieval cues and keys serve as memory traces, with the attention weight between a query-key pair reflecting the degree of match between the cue and the stored trace. This parallels the content-addressed, associative retrieval that characterizes human memory, where memories are accessed through contextual cues rather than through fixed addresses. Research by \citet{ryskin2025attention} has explored these parallels in detail, asking whether attention provides a plausible cognitive model of human memory retrieval and what constraints this imposes on the nature of the underlying memory representations.

Attention weights have been shown to track relational complexity in ways that parallel the working memory demands of human cognitive tasks. When transformers process sentences with complex relational structures (e.g., nested relative clauses or multiple interacting entities), the pattern of attention weights reflects the complexity of the relational processing required, serving as a computational analogue of working memory load. Neuroimaging studies have identified structural correspondences between transformer components and brain regions: positional encoding representations align with visual cortex representations of spatial structure, while attention head activation patterns map onto the frontoparietal and default-mode networks that are associated with attentional control and memory retrieval in human neuroimaging studies \citep{nikolaev2025aligning}.

\subsubsection{TransformerFAM and Feedback Attention}

\citet{irie2024transformerfam} introduced TransformerFAM (Feedback Attention Memory), which explicitly models feedback loops within the attention mechanism to simulate working memory through iterative refinement. Standard transformer attention is a feedforward process: information flows from input to output in a single pass through each layer. TransformerFAM adds recurrent feedback connections that allow the output of attention in one pass to influence the attention computation in subsequent passes, creating a form of iterative refinement that maintains and updates internal representations over multiple processing cycles. This recurrent processing within attention layers provides closer alignment with cognitive theories of active maintenance, where working memory contents are not passively stored but actively refreshed and updated through recurrent neural activity.

\paragraph{Hebbian Learning and Attention.} Research by \citet{kozachkov2022hebbian} has demonstrated that short-term Hebbian learning---a form of synaptic plasticity in which connections between simultaneously active neurons are strengthened---can implement computations that are mathematically equivalent to transformer-like attention. This finding suggests that the computational principle underlying both biological and artificial attention may be fundamentally similar: both implement a form of content-addressed associative retrieval through temporary modifications of connection strengths. If this interpretation is correct, it provides a principled biological foundation for the use of attention as a working memory mechanism in artificial systems, and suggests that the capacity limitations observed in both human working memory and transformer attention may reflect deep computational constraints inherent in content-addressed associative retrieval.

\paragraph{Disanalogies.} Important disanalogies between transformer attention and human memory retrieval must also be acknowledged. Transformer attention requires the quadratic computation of all pairwise similarities between queries and keys, a process with no clear biological counterpart---human memory retrieval does not require exhaustive comparison of a retrieval cue against all stored traces. Human memory retrieval is content-addressed and probabilistic, with retrieval success depending on encoding conditions, retrieval cue specificity, and interference from competing memories, while transformer attention is deterministic and exhaustive within its context window. Human memory exhibits encoding specificity effects, generation effects, and testing effects that have no clear analogues in standard attention mechanisms. These disanalogies define opportunities for architectural innovation: incorporating stochastic retrieval, encoding-specificity constraints, or retrieval-practice effects into attention mechanisms could yield systems that more closely approximate the functional properties of human memory.

\subsection{Episodic versus Semantic Memory in LLMs}

The Tulving distinction between episodic and semantic memory \citep{tulving1972episodic, tulving1985memory} has profound implications for LLM design that extend beyond the theoretical framework reviewed above. This subsection examines the specific challenges that LLMs face with episodic memory and the solutions that have been proposed to address them.

\subsubsection{The Episodic Memory Deficit in Standard LLMs}

Standard LLMs encode semantic memory effectively: the vast knowledge base acquired during pre-training captures general truths, factual relationships, and linguistic patterns in the distributed representations of the model's parameters. However, they are fundamentally deficient in episodic memory. They cannot remember specific interactions without external augmentation. They cannot track the temporal context of information---when something was said, in what order events occurred, or how beliefs have changed over time. They cannot maintain spatiotemporal grounding across sessions, creating the illusion of a continuous agent identity only within the boundaries of a single context window \citep{ning2025episodic, sun2025episodic}.

Recent empirical work has systematically characterized these deficits. State-of-the-art models including GPT-4, Claude, and Llama 3.1 struggle with episodic memory tasks involving multiple related events, complex spatiotemporal relationships, and person-specific memories \citep{ning2025episodic}. When multiple related events are presented, models have difficulty tracking causal chains and maintaining consistent temporal ordering. When person-specific information is involved, models frequently confuse who said what, merging information from different sources into an undifferentiated semantic representation. These failures are not merely limitations of context window size: they reflect the absence of dedicated episodic memory mechanisms that would maintain the temporal, contextual, and source-specific information that distinguishes episodic from semantic memory.

\subsubsection{Architectural Solutions for Episodic Memory}

Several architectural approaches have been proposed to address the episodic memory deficit. External episodic memory systems, exemplified by EM-LLM \citep{hu2025emllm} and the memory stream of Generative Agents \citep{park2023generative}, provide explicit storage for temporally grounded events with associated metadata (timestamps, source information, importance scores). These systems implement computational versions of the episodic memory system described by Tulving, with retrieval mechanisms that weight recency, relevance, and importance in ways that parallel the cue-dependent, context-sensitive retrieval that characterizes human episodic memory.

Explicit temporal tagging and event segmentation provide another approach. EM-LLM \citep{hu2025emllm} uses Bayesian surprise combined with graph-theoretic boundary refinement to segment continuous experience into discrete events, paralleling the event segmentation processes that cognitive scientists have identified in human perception and memory. This segmentation enables temporally contiguous retrieval: when a relevant event is identified, surrounding events can also be retrieved, capturing the temporal context that is essential for episodic memory function. Structured retrieval mechanisms that combine multiple signals---recency, relevance, and importance---outperform any single retrieval signal, as demonstrated by both Generative Agents and EM-LLM. This multi-signal approach reflects the cognitive science finding that human memory retrieval is influenced by multiple factors simultaneously, including recency, emotional significance, and contextual similarity.

The integration of episodic and semantic memory represents the most ambitious research direction. The ideal architecture would support both the storage and retrieval of specific episodes (who said what, when, in what context) and the gradual consolidation of episodic information into semantic knowledge (general patterns, user preferences, domain expertise). This integration mirrors the biological process by which episodic memories are gradually consolidated into semantic knowledge through hippocampal-cortical replay, and it requires mechanisms for identifying when episodic information has been sufficiently consolidated to be represented more efficiently in semantic form.

\subsection{Continual Learning, Catastrophic Forgetting, and Long-Term Memory}

The challenge of maintaining stable long-term knowledge while acquiring new information connects directly to the CLS framework reviewed above. This subsection examines catastrophic forgetting as the computational manifestation of the interference problem that motivated the evolution of complementary learning systems in biological brains, and reviews the strategies that have been developed to mitigate it.

\subsubsection{The Nature and Severity of Catastrophic Forgetting}

When neural networks are trained on new tasks or data, they often catastrophically forget previously learned information \citep{luo2023empirical}. This phenomenon, known as catastrophic forgetting or catastrophic interference, arises because gradient-based learning updates parameters to optimize performance on the current training objective, potentially overwriting the parameter configurations that were optimal for previous objectives. The mathematical basis of the problem is straightforward: the update rule $w_{\text{new}} = w_{\text{old}} - \eta \nabla L(w_{\text{old}})$ directly modifies the weights that encoded previous knowledge, and there is no intrinsic mechanism to preserve the information content of the original configuration while adapting to new requirements.

In large language models, catastrophic forgetting manifests through several interrelated mechanisms. Gradient interference in attention weights disrupts the carefully calibrated attention patterns that enable effective processing of previously learned tasks. Representational drift in intermediate layers shifts the internal representations away from configurations that were effective for prior knowledge. Loss landscape flattening around prior task minima reduces the model's ability to maintain sharp, discriminative representations for old tasks while learning new ones. Empirical studies have documented the severity of these effects: \citet{luo2023empirical} found that the severity of catastrophic forgetting increases with model size in the 1B to 7B parameter range, with fifteen to twenty-three percent of lower-layer attention heads showing severe disruption during continual fine-tuning. This scale-dependent effect is particularly concerning because it suggests that the problem will not be solved merely by scaling models to larger sizes.

The relationship between catastrophic forgetting and memory consolidation is fundamental. Catastrophic forgetting is, in essence, the computational manifestation of the problem that motivated the evolution of complementary learning systems in biological brains. LLMs lack the biological mechanisms---separate fast-learning and slow-learning systems, sleep-dependent consolidation, interleaved replay---that prevent interference in human memory. This absence makes them vulnerable to forgetting in ways that humans are not, and addressing this vulnerability requires implementing computational analogues of the consolidation mechanisms that biological evolution developed to solve the same problem.

\subsubsection{Parameter-Efficient Fine-Tuning and LoRA}

Low-Rank Adaptation (LoRA) \citep{hu2022lora} represents one approach to reducing catastrophic forgetting by limiting the magnitude and scope of parameter changes during fine-tuning. LoRA freezes the pre-trained model weights and injects small, trainable low-rank matrices into the model's layers. The parameter update is decomposed as $\Delta W = BA$, where $B$ is an $r \times d$ matrix and $A$ is a $d \times r$ matrix, with the rank $r$ typically set to 8, 16, or 32. This decomposition reduces the number of trainable parameters by approximately 10,000-fold compared to full fine-tuning and reduces GPU memory requirements by approximately 3-fold, enabling fine-tuning of models with tens of billions of parameters on a single GPU.

However, recent research has revealed that LoRA alone does not prevent catastrophic forgetting. While LoRA reduces forgetting compared to full fine-tuning, it still suffers from significant knowledge degradation: empirical studies have shown that standard LoRA fine-tuning can result in seventy-one percent loss of original task knowledge, compared to eighty-nine percent for full fine-tuning. Advanced variants such as OPLoRA (Orthogonal Projection LoRA), which constrains updates to be orthogonal to the most important directions of the original weight matrices, and CL-LoRA (Continual Learning LoRA), which maintains shared adapter layers with task-specific gating mechanisms, have achieved substantially better results, reducing forgetting to approximately five to forty-five percent depending on the approach and evaluation conditions.

The relationship between LoRA and biological memory consolidation is instructive. LoRA acts as a ``minimal perturbation'' to parametric memory---a small, structured modification that preserves most of the original knowledge while allowing targeted adaptation. This parallels the biological principle that synaptic modifications during learning should be minimal and targeted, changing only the connections necessary for new learning while preserving the broader network structure. However, LoRA lacks the temporal dynamics of biological consolidation: it applies all modifications simultaneously rather than gradually integrating new knowledge through interleaved replay, and it has no mechanism for identifying and protecting the most important parameters for previously learned tasks.

\subsubsection{Experience Replay}

Experience replay maintains buffers of previous training examples and interleaves them with new training data, directly paralleling hippocampal replay in biological memory consolidation. The connection to \citet{mcclelland1995complementary} is explicit: replay prevents catastrophic interference by ensuring that the neocortical (parametric) learning system is exposed to a mixture of old and new experiences, rather than being trained exclusively on new data that would overwrite old knowledge. Sparse memory fine-tuning techniques have demonstrated that storing and replaying only two percent of original training data during new task learning can reduce catastrophic forgetting from seventy-one percent (with LoRA alone) to eleven percent, while maintaining equivalent new task learning. This dramatic improvement with minimal additional storage supports the CLS prediction that interleaved replay is the most effective mechanism for preventing catastrophic interference.

\subsubsection{Elastic Weight Consolidation}

Elastic Weight Consolidation (EWC) \citep{kirkpatrick2017overcoming} implements a regularization approach inspired by synaptic consolidation in biological neural networks. EWC computes the Fisher information matrix for each parameter after learning a task, identifying which parameters are most important for that task's performance. During subsequent learning, a quadratic penalty prevents these important parameters from changing significantly:
\begin{equation}
L_{\text{total}}(w) = L_{\text{new}}(w) + \frac{\lambda}{2} \sum_i F_i (w_i - w_i^*)^2
\end{equation}
where $w_i^*$ are the optimal parameters for the previous task, $F_i$ is the Fisher information for parameter $i$, and $\lambda$ controls the strength of the consolidation penalty. This approach is ``elastic'' in the sense that it allows some parameter movement (unlike rigid freezing) but penalizes movement in proportion to the parameter's importance for previously learned tasks.

The biological analogue of EWC is synaptic tagging and capture, a mechanism by which recently active synapses are marked (tagged) for consolidation, making them resistant to modification during subsequent learning. Both EWC and synaptic tagging implement the same principle: not all connections are equally important, and consolidation should selectively protect the most critical ones while leaving others available for new learning. The limitation of EWC is that the Fisher information approximation becomes increasingly inaccurate as the number of sequential tasks grows, and the cumulative regularization can eventually prevent any further learning.

\subsubsection{Dynamic Architectures}

Dynamic architectures, such as progressive neural networks, address catastrophic forgetting by adding new capacity for new tasks without modifying existing parameters. New lateral connections allow the new capacity to access features learned by previous modules, enabling knowledge transfer without interference. This approach has a biological analogue in neurogenesis---the creation of new neurons in the hippocampus throughout life---which has been proposed as a mechanism for maintaining the ability to form new memories without disrupting existing ones. While dynamic architectures eliminate catastrophic forgetting entirely, they incur growing computational and memory costs as the number of tasks increases, making them impractical for systems that must learn many sequential tasks.

\subsection{Sleep-Inspired Consolidation for Artificial Neural Networks}

The biological mechanisms of sleep-dependent consolidation reviewed above have inspired a nascent but potentially transformative research direction: the development of sleep-like offline processing phases for artificial neural networks.

\subsubsection{NeuroDream: Explicit Dream Phases}

\citet{tutuncuoglu2024neurodream} introduced NeuroDream, which adds explicit ``dream phases'' to neural network training. During dream phases, the model disconnects from input data and engages in internal simulations: it samples from learned latent representations, generates synthetic experiences using a learned dynamics model, and consolidates knowledge through Hebbian-like updates that strengthen important patterns without re-exposure to raw training data. The dream phase comprises four key components. A latent embedding space provides compressed representations of past experiences. A dynamics model learns temporal patterns and predicts state transitions. A dream generator samples from the latent space to produce synthetic experiences. A consolidation objective maximizes the diversity and quality of dream experiences while strengthening generalizable patterns.

The results of NeuroDream experiments demonstrate meaningful reductions in catastrophic forgetting---from forty percent forgetting without dream phases to fifteen percent with them---at a cost of approximately twenty percent additional training time. These results, while preliminary, suggest that offline consolidation phases can address the fundamental limitation of online-only learning in artificial systems. The mechanism parallels biological sleep consolidation in several respects: the disconnection from external input prevents interference between consolidation and new learning, the replay of stored representations provides the interleaved exposure that the CLS framework identifies as essential for preventing catastrophic interference, and the abstraction process that operates over replayed experiences extracts generalizable patterns from specific episodes.

\subsubsection{Sleep Replay Consolidation}

Research on sleep-like unsupervised replay has demonstrated that interleaving periods of offline replay between active learning phases can preserve task knowledge to a remarkable degree. In controlled experiments, networks trained on two sequential tasks showed dramatic forgetting of the first task without sleep replay (accuracy dropping from ninety-two percent to eighteen percent), but maintained near-original performance with sleep replay (eighty-six percent, compared to ninety-one percent for the original task and eighty-three percent for the new task). The sleep replay mechanism involves disconnecting from external data, reactivating stored memory traces through pattern replay, applying Hebbian consolidation to strengthen important connections, and reorganizing representations for improved abstraction. These results align with the biological finding that sleep enhances memory consolidation by twenty-five to forty percent compared to equivalent periods of wakefulness.

\subsubsection{Computational Models of Hippocampal-Neocortical Consolidation}

\citet{gonzalez2022model} developed computational models that simulate the autonomous interactions between hippocampus and neocortex during sleep-dependent consolidation. These models demonstrate that offline learning through simulated replay can build structured representations via interleaved replay, with different simulated sleep stages serving complementary consolidation roles. NREM-like phases, characterized by slow oscillatory dynamics, support the transfer of specific episodic information from hippocampal to cortical representations. REM-like phases, characterized by theta-frequency dynamics, support the integration of transferred information into existing semantic structures and the extraction of abstract patterns. The alternation between these phases---mirroring the ultradian cycling of NREM and REM sleep in biological systems---produces more robust consolidation than either phase alone, suggesting that the temporal organization of consolidation processes is as important as the processes themselves.

These sleep-inspired approaches represent a research direction that could fundamentally change how LLMs are trained and deployed. If offline consolidation phases become standard in LLM training pipelines, they could address the fundamental limitation of catastrophic forgetting while enabling more human-like knowledge integration. The key challenge is computational: current LLMs require enormous resources for training, and adding offline consolidation phases would increase training costs. However, if consolidation phases reduce the total amount of training data required (by enabling more efficient knowledge integration) or improve the quality of the resulting models (by producing more robust, less interference-prone representations), the additional cost may be justified.

\subsection{Infinite Context and Unbounded Memory Architectures}

The desire to overcome the fundamental capacity limitations of the context window has motivated several architectures that aim to provide effectively unbounded context processing while maintaining computational feasibility.

\subsubsection{Infini-Attention}

\citet{munkhdalai2024leave} introduced Infini-Attention, which augments standard attention with a compressive memory layer that enables processing of contexts exceeding one million tokens with 114-fold less memory than standard attention. The architecture maintains a compressed memory state alongside the standard attention computation, allowing the model to retain information from the entire preceding context in a fixed-size compressed representation while performing full attention within a local window. The system demonstrated the ability to retrieve specific information (such as randomly inserted numbers) from million-token contexts and to summarize 500,000-token books effectively, performance levels that are impossible for standard attention architectures at these scales.

Infini-Attention's compressive memory layer functions as a form of long-term memory that operates alongside the working memory of the local attention window. The compression process that maps detailed local representations into the fixed-size memory state parallels the consolidation processes that transform detailed episodic memories into more compact semantic representations in biological memory. The challenge is in the quality of compression: information that is compressed into the memory state may lose the fine-grained detail that would be necessary for certain retrieval tasks, analogous to the loss of episodic detail that occurs as biological memories are consolidated from hippocampal to cortical representations.

\subsubsection{EM-LLM: Human-Inspired Episodic Memory for Infinite Context}

\citet{hu2025emllm} developed EM-LLM, which integrates human-inspired episodic memory and event cognition principles to enable unbounded context processing without fine-tuning. EM-LLM's approach is grounded in the cognitive science of event segmentation, which has established that humans naturally segment continuous experience into discrete events at points of significant change in the ongoing activity.

EM-LLM implements event segmentation using a combination of Bayesian surprise (detecting points where the incoming information significantly deviates from predictions based on the preceding context) and graph-theoretic boundary refinement (adjusting event boundaries to maximize within-event coherence while minimizing between-event similarity). This produces a segmentation of the input text into coherent episodic units that can be individually stored, indexed, and retrieved. Retrieval operates in two stages: similarity-based retrieval identifies the most relevant events for a given query, and temporally contiguous retrieval expands the retrieved context to include events adjacent to the retrieved events, capturing temporal context.

EM-LLM demonstrated remarkable performance: it surpassed full-context models on most evaluation tasks and successfully retrieved relevant information across ten million tokens without any fine-tuning. The system's ability to operate at this scale while maintaining retrieval quality demonstrates that human-inspired event segmentation and episodic memory mechanisms can provide practical solutions to the context scaling challenge. The event segmentation approach is particularly well-suited to the structure of many real-world LLM applications, where the input naturally decomposes into discrete episodes (conversation turns, document sections, task steps) that can be individually stored and retrieved.

\subsubsection{Hierarchical Memory Transformer}

The Hierarchical Memory Transformer (HMT), presented at NAACL 2025, implements a three-tier memory hierarchy within the transformer architecture itself. The architecture comprises a long-term memory tier that stores key-value embeddings from early input positions with sparse attention, a short-term working memory tier that maintains full self-attention within a sliding local window, and a current focus tier that generates predictions by attending across all memory tiers. This hierarchy enables processing of sequences up to 32,000 tokens with 2.5 to 116-fold speedup compared to full attention, while maintaining generation quality comparable to full attention on short contexts and achieving superior quality on long contexts. The efficiency gains arise from replacing the quadratic cost of full attention with the combination of bounded quadratic cost within the local window and linear-cost sparse attention to the long-term memory tier.

\subsubsection{Hierarchical Context Caching}

Strata, a hierarchical context caching system for LLM serving, applies memory hierarchy principles at the infrastructure level. In production environments where LLM servers must maintain key-value caches for ongoing conversations, Strata implements a three-tier storage hierarchy: GPU VRAM (hot cache) for the most recent tokens with full precision and immediate access, CPU RAM (warm cache) for mid-range history in compressed form, and disk/NVMe storage (cold cache) for older context with high compression. A least-recently-used eviction policy promotes frequently accessed entries to faster tiers while demoting unused entries to slower, cheaper storage. This infrastructure-level memory hierarchy mirrors the cognitive memory hierarchy at a systems engineering level, demonstrating that the principles of hierarchical memory organization apply not only to the algorithms that process information but also to the systems that store and serve it.

\subsection{Comparative Analysis: Human versus LLM Memory}

A systematic comparison of human and LLM memory systems reveals both striking parallels and fundamental differences that inform the design of future architectures. The following analysis synthesizes findings from comparative studies \citep{cowan2024judgments, arrieta2024llm_cognitive} with the theoretical frameworks reviewed above.

\begin{longtable}{L{3.5cm}L{4.5cm}L{5cm}}
\toprule
\textbf{Property} & \textbf{Human Memory} & \textbf{LLM Memory} \\
\midrule
\endfirsthead
\toprule
\textbf{Property} & \textbf{Human Memory} & \textbf{LLM Memory} \\
\midrule
\endhead
Sensory Memory & Large capacity; 100--500ms retention; acts as input gate & Token embeddings at API input; immediate processing without persistence \\
\midrule
Working Memory & 4--7 chunks; seconds to minutes; active manipulation & Context window (512--128K+ tokens); no inter-session persistence; quadratic attention cost \\
\midrule
Long-Term Memory & Unlimited capacity; lifetime; dynamic, reconstructive; susceptible to distortion & Parametric (fixed post-training) + non-parametric (external databases, RAG); static until re-training \\
\midrule
Memory Formation & Encoding via attention and elaboration; hippocampal binding; consolidation over hours to years & Training (parametric) or direct storage in external memory (non-parametric); no consolidation process \\
\midrule
Consolidation & Sleep-dependent; hippocampal replay; gradual neocortical integration; NREM/REM complementarity & Absent in standard LLMs; experimental in NeuroDream and sleep replay consolidation research \\
\midrule
Forgetting & Exponential decay (Ebbinghaus); adaptive and importance-weighted; interference-based; retrieval failure & No intrinsic forgetting mechanism; catastrophic interference during fine-tuning; context window truncation \\
\midrule
Episodic Memory & Temporally grounded events; autonoetic consciousness; rich contextual detail; reconstructive & External systems (EM-LLM, Generative Agents); context-dependent; no phenomenal experience \\
\midrule
Semantic Memory & General knowledge; noetic consciousness; conceptually stable across contexts & Parametric knowledge; static until re-training; context-dependent activation patterns \\
\midrule
Retrieval & Cue-based; reconstructive; emotion-influenced; encoding specificity; testing effect & Attention-based (parametric) or similarity search (non-parametric); deterministic; prompt-sensitive \\
\midrule
Statefulness & Continuous, persistent identity; autobiographical continuity & Stateless by default; each inference independent; memory systems approximate continuity \\
\midrule
Selectivity & High; emotional significance modulates encoding via amygdala-hippocampal interaction & Low; all tokens receive equal computational resources; no intrinsic importance weighting \\
\midrule
Metacognition & Calibrated judgments of learning; accurate confidence monitoring & Poorly calibrated; confidence-recall correlation near zero in empirical studies \\
\midrule
Adaptability & Context-sensitive; continuous learning; personal biases shape recall & Static parametric knowledge; adaptable via external memory or fine-tuning \\
\midrule
Knowledge Update & Continuous; synaptic plasticity; no catastrophic forgetting under normal conditions & Requires re-training or PEFT; catastrophic forgetting during continual learning \\
\midrule
Computational Cost & Biological efficiency; parallel processing; approximately 20 watts & Quadratic attention complexity; high energy consumption; significant memory bottlenecks \\
\bottomrule
\caption{Comprehensive comparison of human memory and LLM memory systems, synthesizing findings across cognitive science and machine learning research.}
\label{tab:human_llm_memory}
\end{longtable}

\subsubsection{Key Insights from the Comparison}

Several insights emerge from this systematic comparison that have direct implications for LLM memory design.

\paragraph{The Statefulness Gap.} The most fundamental difference between human and LLM memory is statefulness. Human memory provides continuous, persistent identity across time---the sense of being the same person who went to sleep last night and woke up this morning. Standard LLMs are fundamentally stateless: each inference begins with a blank slate (or, at most, a context window filled with the current prompt), and there is no intrinsic mechanism for carrying information from one session to the next. Memory-augmented systems such as MemGPT and Generative Agents approximate statefulness by maintaining external stores that persist across sessions, but this approximation is qualitatively different from the continuous, embodied persistence of human identity. The statefulness gap is particularly significant for applications that require long-term agent deployment, where the agent must maintain consistent identity, accumulated knowledge, and relationship history across many interactions.

\paragraph{The Consolidation Gap.} Sleep-dependent consolidation is central to human memory but entirely absent in standard LLMs. The consequences of this gap are severe: without consolidation, LLMs cannot gradually integrate new knowledge with existing representations, cannot extract abstract patterns from specific experiences, and cannot selectively strengthen important memories while allowing unimportant ones to fade. The experimental work on NeuroDream, sleep replay consolidation, and computational models of hippocampal-neocortical interaction represents the beginning of an effort to close this gap, but these approaches are far from practical deployment in large-scale LLM systems.

\paragraph{The Metacognition Gap.} Empirical research has demonstrated that LLMs exhibit near-zero correlation between their confidence judgments and their actual recall accuracy, in stark contrast to humans, who show strong correlations between judgments of learning and subsequent recall. This metacognitive deficit means that LLMs cannot reliably assess whether they ``know'' something or are confabulating, a limitation that has direct implications for the reliability of memory-augmented systems. If an LLM cannot distinguish between genuine memory retrieval and hallucinated information, the value of its memory system is fundamentally compromised.

\paragraph{The Selectivity Gap.} Human memory is highly selective, preferentially encoding emotionally significant, goal-relevant, and novel information through amygdala-hippocampal interactions. LLMs, by contrast, treat all information with approximately equal computational resources, lacking any intrinsic mechanism for importance-weighted encoding. Memory architectures that implement importance scoring (such as Generative Agents and MemoryBank) partially address this gap, but their importance ratings are assigned post hoc by the language model rather than being integrated into the encoding process itself.

\paragraph{Surprising Convergence.} Despite these profound differences, certain properties show surprising convergence. The effective working memory capacity of transformer attention (approximately four to seven actively considered items) is strikingly similar to the human working memory capacity measured by \citet{cowan2001magical}. Both humans and LLMs exhibit serial position effects (primacy and recency). Both generate ``false memories'' (human memory distortion paralleling LLM hallucination). These convergences suggest that certain properties of memory systems may emerge from computational constraints that are shared between biological and artificial information processing, rather than being specific to either substrate.

\subsection{Integrated Multi-Tier Memory Systems}

Building on the individual architectures reviewed above, recent work has begun to develop integrated multi-tier memory systems that combine multiple cognitive-science-inspired principles into unified architectures.

\subsubsection{Hierarchical Memory for Agents}

Recent work on hierarchical memory for LLM agents has proposed four-layer memory structures that organize information by level of abstraction. A domain layer provides broad semantic categorization (work, personal, health). A category layer provides finer-grained sub-domain grouping within each domain. A memory trace layer stores abstracted summaries of key information. An episode layer preserves complete, detailed interaction records. This hierarchical organization enables top-down search strategies with logarithmic rather than linear complexity: a query about a work deadline can navigate from the ``work'' domain to the ``projects'' category to relevant memory traces without exhaustively scanning all stored memories.

This hierarchical organization mirrors the semantic network models of human memory that posit hierarchical categorical structures for organizing knowledge, and it implements a form of the chunking principle identified by \citet{miller1956magical}: by organizing memories into meaningful categories, the system effectively increases the amount of information that can be efficiently accessed despite limited processing resources.

\subsubsection{Cognitive Workspace Theory Applied to LLMs}

Inspired by Baars's Global Workspace Theory of consciousness, the cognitive workspace approach treats the LLM's context window as a limited-capacity workspace to which multiple cognitive modules compete for access. The workspace holds the current task goal, recent observations, retrieved relevant knowledge, and active reasoning traces, with a capacity of approximately 1,000 to 4,000 tokens. Multiple cognitive modules---reading, planning, memory, and execution modules---operate in parallel, with each module contributing information to the workspace and using workspace contents to guide its own processing. The workspace functions as a broadcast mechanism: information placed in the workspace becomes available to all modules simultaneously, enabling implicit coordination without explicit communication between modules.

This approach transcends traditional retrieval-augmented generation by treating memory not as a static store to be queried but as an active computational resource to be managed. The workspace metaphor shifts the focus from ``what is stored'' to ``what is currently in the spotlight of attention,'' a shift that aligns with the cognitive science understanding that working memory is not merely a passive buffer but an active processing system that shapes and transforms the information it contains.

\subsubsection{Neuro-Symbolic Integration}

The integration of symbolic knowledge representations (knowledge graphs, production rules, logical constraints) with neural language models represents a hybrid approach that leverages the complementary strengths of both paradigms. Symbolic representations provide interpretability, precise relational reasoning, and guaranteed consistency, while neural representations provide flexibility, generalization, and natural language understanding. Hybrid architectures that combine a symbolic layer for structured knowledge and constraint enforcement with a neural layer for language understanding and flexible reasoning can achieve properties that neither layer provides alone: grounded generation that respects factual constraints, interpretable decision-making that can be audited and explained, and flexible reasoning that can handle novel situations not covered by explicit rules.

The connection to cognitive architecture is direct: the symbolic layer corresponds to the declarative and procedural memory systems of architectures like ACT-R and SOAR, while the neural layer corresponds to the pattern-recognition and association capabilities that connectionist models have long provided. The integration of these two layers---allowing symbolic knowledge to constrain neural generation while neural processing discovers new patterns to be represented symbolically---mirrors the interaction between explicit and implicit memory systems in human cognition.

\subsection{Synthesis: Design Principles for Human-Inspired LLM Memory}

The convergence of cognitive science and LLM research reviewed in this section yields a coherent set of design principles for the construction of effective, scalable, and human-like memory systems for large language models.

\subsubsection{The Multi-Tier Imperative}

Effective memory systems require, at minimum, three tiers of storage: a working memory tier (the context window) for immediate processing, an episodic memory tier (external storage with temporal and contextual metadata) for specific events and experiences, and a semantic memory tier (parametric knowledge and structured knowledge bases) for generalized facts and concepts. Single-tier designs---whether based on context window expansion, pure parametric knowledge, or undifferentiated external storage---are fundamentally limited because they cannot simultaneously satisfy the competing requirements of rapid access, temporal grounding, and stable, generalized knowledge representation. The biological evidence reviewed in this section strongly supports the multi-tier approach: evolution produced multiple memory systems with distinct temporal and representational properties precisely because no single system can meet all the requirements of intelligent behavior.

\subsubsection{Selective Consolidation and Principled Forgetting}

Not all information should persist indefinitely. Importance-weighted storage, as implemented by Generative Agents and MemoryBank, and temporal decay, as inspired by the Ebbinghaus forgetting curve, are essential mechanisms for preventing unbounded memory growth while preserving the information that is most likely to be relevant in the future. The principle of selective consolidation---that memory management should actively decide what to retain, what to compress, what to abstract, and what to forget---is supported by both the cognitive science of adaptive forgetting and the information-theoretic principle that optimal information retention under resource constraints requires selective rather than exhaustive storage. Forgetting is not a deficiency to be eliminated but a feature to be designed, and effective LLM memory systems must implement principled forgetting mechanisms that balance the cost of storage and retrieval against the value of retained information.

\subsubsection{Temporal Awareness}

Timestamps, decay functions, and temporal ordering are critical for episodic memory function. Without temporal awareness, memory systems cannot distinguish between recent and remote information, cannot track the evolution of beliefs over time, and cannot implement the recency-weighted retrieval that is essential for contextually appropriate behavior. Systems such as RecallM, MemoryBank, and EM-LLM demonstrate the importance of temporal information for effective memory management, and the cognitive science evidence on the role of temporal context in human episodic memory strongly supports the integration of temporal awareness into LLM memory architectures.

\subsubsection{Multi-Signal Retrieval}

Combining recency, relevance, and importance in retrieval, as demonstrated by Generative Agents and EM-LLM, outperforms any single retrieval signal. This multi-signal approach reflects the cognitive science finding that human memory retrieval is influenced by multiple factors simultaneously, and it provides robustness against the failure of any individual signal. A memory that is old but highly important and relevant to the current query should be retrieved, just as a memory that is recent but unimportant and irrelevant should not be. Multi-signal retrieval enables this kind of flexible, context-sensitive access to stored information.

\subsubsection{Dynamic Belief Revision}

Memory systems must support the revision of stored knowledge when new information contradicts previous beliefs. Systems such as RecallM and RET-LLM demonstrate that knowledge representation formats that support explicit belief revision (graph databases, semantic triplets) are more effective for dynamic knowledge management than formats that only support accumulation (vector databases, unstructured text). The ability to update, revise, and correct stored knowledge is essential for agents operating in dynamic environments where facts change over time, and it parallels the human capacity for belief revision in the face of new evidence.

\subsubsection{Reflection and Abstraction}

Higher-order processing that generates summaries, insights, and abstractions from accumulated experience, as demonstrated by Generative Agents and Think-in-Memory, dramatically improves agent performance. Reflection transforms raw episodic data into actionable knowledge, reducing the volume of information that must be searched during retrieval while preserving and enhancing its informational value. The cognitive science analogue is the consolidation process that transforms specific episodic memories into generalized semantic knowledge, and the evidence reviewed in this section suggests that this transformation is not merely a storage optimization but a fundamental cognitive operation that enables higher-level reasoning, planning, and decision-making.

\subsubsection{Controlled Memory Integration}

Memory controllers that evaluate the relevance of retrieved information before injecting it into the context, as implemented by SCM and MemGPT, are necessary for preventing noise accumulation during long-term operation. The injection of irrelevant or outdated information into the context window can degrade performance rather than improve it, and intelligent filtering at the point of integration is as important as effective storage and retrieval. This principle parallels the attentional gating function of the central executive in Baddeley's working memory model, which selectively admits information to working memory based on its relevance to current processing goals.

\subsubsection{Event Segmentation}

Chunking continuous experience into meaningful episodes, as implemented by EM-LLM and motivated by CoALA, enables efficient retrieval by providing natural units of storage and search. Event segmentation mirrors the event perception processes that cognitive scientists have identified in human experience, where continuous sensory input is automatically segmented into discrete events at points of significant change. By storing memories as coherent episodic units rather than as undifferentiated streams of tokens, memory systems can retrieve contextually complete episodes that provide the temporal and situational context necessary for informed decision-making.

\subsection{Future Directions and Open Problems}

The synthesis of cognitive science and LLM memory design reviewed in this section points to several critical open problems and research directions that will shape the next generation of memory-augmented language systems.

\paragraph{Biological Plausibility.} The gap between biologically-inspired design principles and their computational implementations remains wide. While systems like MemoryBank implement simplified versions of the Ebbinghaus forgetting curve and Generative Agents implement a form of reflection that parallels consolidation, no existing system implements the full complexity of biological memory consolidation, including the coordinated oscillatory dynamics of sleep, the complementary roles of NREM and REM sleep, or the pattern separation and completion mechanisms of hippocampal circuits. Closing this gap will require closer collaboration between cognitive neuroscientists and machine learning researchers.

\paragraph{Scalable Consolidation.} If offline consolidation phases are to become practical for large language models, significant advances in computational efficiency will be required. Current sleep-inspired approaches add twenty to fifty percent to training time, which is prohibitive at the scale of frontier LLMs. Developing consolidation mechanisms that operate efficiently within existing training pipelines, or that can be applied selectively to the most important information, is an essential research challenge.

\paragraph{Unified Memory Architectures.} No existing system seamlessly integrates parametric semantic memory, external episodic memory, and procedural learning within a single coherent framework. CoALA provides the conceptual vocabulary for such integration, but practical implementations remain fragmented across different systems with different storage formats, retrieval mechanisms, and update procedures. Developing unified architectures that support the full range of memory operations---encoding, consolidation, retrieval, updating, and forgetting---across multiple memory types is perhaps the most important challenge for the field.

\paragraph{Metacognitive Monitoring.} The near-zero correlation between LLM confidence and recall accuracy identified in empirical studies represents a fundamental limitation for memory-augmented systems. If an agent cannot assess whether a retrieved memory is accurate, the reliability of its memory-informed decisions is fundamentally compromised. Developing mechanisms for metacognitive monitoring---enabling LLMs to assess the reliability of their own memories and retrieval processes---is essential for trustworthy deployment of memory-augmented agents.

\paragraph{Episodic Richness.} Current external episodic memory systems store relatively impoverished representations of events compared to the rich, multimodal, emotionally tagged episodic memories that humans maintain. Capturing the full richness of episodic experience---including emotional valence, sensory detail, spatial context, and social dynamics---in a computationally tractable format remains an open challenge.

These principles and open problems suggest that the most effective context management systems will not merely extend context windows or add external storage. They will implement biologically-inspired hierarchical memory architectures with principled consolidation, retrieval, and forgetting mechanisms that draw on the deep insights that cognitive science has accumulated over more than a century of research on the nature and function of human memory.
